% FORM-FB-01: Fichas bibliográficas
% Compilar: pdflatex fichas -> biber fichas -> pdflatex fichas -> pdflatex fichas
\documentclass[11pt,a4paper]{article}

\usepackage[T1]{fontenc}
\usepackage[utf8]{inputenc}
\usepackage{lmodern}
\usepackage[american,spanish,es-noquoting]{babel}
\usepackage{csquotes}
\usepackage[margin=2.5cm]{geometry}
\usepackage{graphicx}
\usepackage{xcolor}
\usepackage{enumitem}
\usepackage[style=ieee,backend=biber,language=american,autolang=other,
  maxnames=6,minnames=1,maxcitenames=6,mincitenames=1]{biblatex}
\addbibresource{referencias.bib}
% Páginas sin separador de miles
\renewcommand*{\mkpagegrouped}[1]{#1}
\usepackage[hidelinks]{hyperref}

% ---------- Datos del documento ----------
\newcommand{\estudiante}{Zambrano Yong Angel Daniel}
\newcommand{\docente}{Ing. Gleiston Cicerón Guerrero Ulloa}
\newcommand{\curso}{Séptimo nivel, paralelo B}
\newcommand{\fecha}{8 de noviembre de 2026}

% ---------- Campos de cada ficha ----------
\newcommand{\campo}[2]{\noindent\textbf{#1:} #2\par\smallskip}
\newcommand{\ficha}[1]{\section*{#1}}
\newcommand{\parte}[1]{\subsection*{#1}}

\setlength{\parindent}{1.5em}

\begin{document}

% =================== CARÁTULA ===================
\begin{titlepage}
  \centering
  \IfFileExists{img/logo-uteq.jpg}{\includegraphics[width=3.5cm]{img/logo-uteq.jpg}\par}{}
  \vspace{1cm}
  {\LARGE\bfseries UNIVERSIDAD TÉCNICA ESTATAL DE QUEVEDO\par}
  \vspace{0.4cm}
  {\large\bfseries Facultad de Ciencias de la Computación\\Carrera de Software\par}
  \vspace{1.2cm}

  {\bfseries Asignatura}\\ Aplicaciones Distribuidas\par\medskip
  {\bfseries Actividad}\\ FORM-FB-01: Fichas bibliográficas sobre fundamentos de sistemas\\
  distribuidos y arquitecturas distribuidas\par\medskip
  {\bfseries Docente}\\ \docente\par\medskip
  {\bfseries Estudiante}\\ \estudiante\par\medskip
  {\bfseries Curso}\\ \curso\par\medskip
  {\bfseries Repositorio}\\
  {\scriptsize\url{https://github.com/Kinetro/QZ-Practicas_Experimentales/tree/main/Semana\%202/FORM-FB-01}\par}

  \vfill
  {\bfseries Año lectivo 2026--2027}\\
  Quevedo -- Los Ríos -- Ecuador\\
  \fecha\par
\end{titlepage}

% =================== PRESENTACIÓN ===================
\section*{Presentación}

Estas fichas reúnen las fuentes que usaré en el informe TA-IND-01, donde analizaré
Apache Cassandra y lo compararé con MongoDB. Escogí tres artículos experimentales que
miden el costo de la consistencia en bases NoSQL, una revisión sobre los compromisos
de la replicación (CAP y PACELC) y otra sobre consenso tolerante a fallos. Todas las
fuentes son de 2021 a 2025.

Cada ficha tiene la referencia en formato IEEE, generada con \texttt{biblatex} desde el archivo
\texttt{.bib}, el DOI, un resumen, un dato concreto con su ubicación y una nota personal.

\clearpage

% =================== FICHA 1 ===================
\ficha{Ficha 1: Niveles de consistencia en bases NoSQL distribuidas en la nube}

\campo{Referencia (IEEE)}{\fullcite{ferreira2025benchmarking}}
\campo{Identificador}{DOI: \href{https://doi.org/10.1109/ACCESS.2025.3558923}{10.1109/ACCESS.2025.3558923}}
\campo{Tipo de fuente}{Artículo de revista indexada (IEEE Access)}

\parte{Resumen}
Es un estudio experimental que mide cuánto rendimiento pierden Cassandra, MongoDB y
Redis cuando se les exige consistencia fuerte. Los trabajos anteriores solían evaluar una
sola base o dejaban fuera las réplicas repartidas en varias regiones. Montaron en Google
Cloud un nodo primario y dos réplicas, cada uno en una región distinta de Estados Unidos,
y corrieron las cargas de YCSB con consistencia débil y fuerte, entre 10 mil y 100 mil
operaciones. Con consistencia fuerte las tres bases procesan muchas menos operaciones por
segundo y tardan bastante más. Redis es la más rápida, pero también la que más cae al
cambiar de nivel.

\parte{Dato concreto}
Promediando todas las cargas, pasar de consistencia débil (ONE) a fuerte (ALL) bajó el
throughput de Cassandra en 64,3\,\% y subió su tiempo de ejecución en 700,71\,\%. En
MongoDB las cifras fueron 60,49\,\% y 719,8\,\%, así que en promedio las dos pagan un precio
parecido.

\campo{Ubicación}{Sección V-B (Performance summary) y Tabla 4, p.~63434.}

\parte{Nota personal}
No esperaba que el costo promedio fuera casi igual en Cassandra y en MongoDB; pensaba
que Cassandra, al estar diseñada para priorizar la disponibilidad, iba a sufrir mucho más.
Voy a usar este dato en la tabla comparativa del TA-IND-01, en la fila de rendimiento con
consistencia fuerte. También me sirve para no exagerar en el informe, porque la diferencia
entre las dos aparece en cargas puntuales, como la D, donde Cassandra cae 75\,\% y
MongoDB 64\,\%.

% =================== FICHA 2 ===================
\ficha{Ficha 2: Impacto de la consistencia en NoSQL geo-replicado}

\campo{Referencia (IEEE)}{\fullcite{ferreira2024impacts}}
\campo{Identificador}{DOI: \href{https://doi.org/10.1186/s13677-024-00716-7}{10.1186/s13677-024-00716-7}}
\campo{Tipo de fuente}{Artículo de revista indexada (Journal of Cloud Computing, Springer)}

\parte{Resumen}
Es un experimento de los mismos autores de la ficha 1 sobre cómo cambia el rendimiento
de Cassandra, MongoDB y Redis según el nivel de consistencia cuando las réplicas están
lejos entre sí. Querían medir el efecto de la distancia entre nodos, que casi no se había
estudiado. Desplegaron las bases en Google Cloud en una sola región y en varias regiones,
y generaron carga con Locust de 1 a 10 usuarios concurrentes, en lecturas y escrituras. En
una región, la consistencia fuerte afecta sobre todo las lecturas de Cassandra. Con varias
regiones todas pierden mucho, y las escrituras de MongoDB son las más golpeadas.

\parte{Dato concreto}
En el despliegue multirregión, las escrituras de Cassandra bajaron de unas 1300 peticiones
por segundo con consistencia débil a unas 25 con consistencia fuerte, una caída del 98\,\%
con 2 usuarios concurrentes. En sus peores casos, MongoDB cayó 94\,\% y Redis 95\,\%.

\campo{Ubicación}{Sección Multi-regions deployment, apartado Throughput in writing
operations, p.~11.}

\parte{Nota personal}
Me llamó la atención que en una sola región Cassandra aguantaba bien las escrituras y que
la caída fuerte apareció recién al separar las réplicas por regiones. Eso me hizo ver que
en el TA-IND-01 tengo que tratar por separado el nivel de consistencia y la ubicación de
los nodos, porque mezclados no se entiende cuál de los dos causa la pérdida. Voy a usar
este dato en la sección sobre replicación geográfica.

% =================== FICHA 3 ===================
\ficha{Ficha 3: Consistencia y compromisos en bases de datos replicadas}

\campo{Referencia (IEEE)}{\fullcite{ahmed2023consistency}}
\campo{Identificador}{DOI: \href{https://doi.org/10.32620/reks.2023.2.14}{10.32620/reks.2023.2.14}}
\campo{Tipo de fuente}{Artículo de revisión en revista de acceso abierto indexada en Scopus
y DOAJ (Radioelectronic and Computer Systems)}

\parte{Resumen}
Es un artículo de revisión sobre los compromisos que aparecen al replicar una base de
datos en varios nodos, que además clasifica trabajos recientes sobre almacenamiento
distribuido. Los autores parten de los teoremas CAP y PACELC, que explican el choque
entre consistencia y disponibilidad o latencia, y suman otros pares, como durabilidad
contra rendimiento o disponibilidad contra costo y consumo de energía. Después agrupan
los artículos revisados por temas usando nubes de palabras y una clasificación propia.
Encuentran que la consistencia y la replicación dominan la investigación, y que crece el
interés por modelos híbridos basados en consistencia causal.

\parte{Dato concreto}
Según PACELC, si hay una partición de red la base elige entre disponibilidad y
consistencia, y si no la hay elige entre latencia y consistencia. Con esa regla, los autores
ubican a Cassandra como PA/EL (prefiere disponibilidad y baja latencia) y a MongoDB como
PA/EC (prefiere disponibilidad ante una partición, pero consistencia cuando la red funciona
normal).

\campo{Ubicación}{Sección 1.2 (PACELC Theorem), p.~172.}

\parte{Nota personal}
Conocía lo básico de CAP, pero no PACELC, y me parece más útil porque explica qué pasa
cuando la red funciona bien, que es casi todo el tiempo. Me sorprendió que el mismo
artículo ponga a MongoDB como CP en la lista de CAP y como PA/EC en la de PACELC; lo voy a
comentar en el informe porque muestra que la etiqueta depende de cómo se configure la base.
Esta ficha la uso en el marco teórico del TA-IND-01, antes de la comparación.

% =================== FICHA 4 ===================
\ficha{Ficha 4: Algoritmos de consenso tolerantes a fallos bizantinos}

\campo{Referencia (IEEE)}{\fullcite{zhang2024byzantine}}
\campo{Identificador}{DOI: \href{https://doi.org/10.1145/3636553}{10.1145/3636553}}
\campo{Tipo de fuente}{Artículo de revista indexada (ACM Computing Surveys)}

\parte{Resumen}
Es una revisión sistemática de algoritmos de consenso tolerantes a fallos bizantinos (BFT),
que permiten que un grupo de servidores se ponga de acuerdo aunque algunos se comporten
de forma arbitraria o maliciosa. Con el auge de blockchain aparecieron tantos algoritmos
nuevos que entenderlos se volvió difícil, porque cada uno parte de supuestos y métricas
distintos. Proponen una taxonomía con tres categorías (más eficientes, más robustos y más
disponibles) y comparan la complejidad de mensajes y de tiempo de cada uno. También
describen cada protocolo paso a paso con diagramas. PBFT aparece como la base de casi
todos, aunque su costo en mensajes crece cuadráticamente.

\parte{Dato concreto}
Para tolerar $f$ fallos benignos, protocolos como Paxos o Raft necesitan quórums de $f+1$
nodos. Con fallos bizantinos el quórum mínimo sube a $2f+1$ dentro de un total de
$n = 3f+1$ nodos. Por ejemplo, aguantar un solo nodo malicioso exige al menos cuatro nodos.

\campo{Ubicación}{Sección 2 (Background), párrafo Byzantine quorums, p.~134:5.}

\parte{Nota personal}
Antes pensaba que tolerar fallos era solo cuestión de tener réplicas, y con este artículo
entendí que el tipo de fallo cambia cuántos nodos hacen falta. Cassandra y MongoDB están
pensadas para fallos benignos, como la caída de un nodo, así que no voy a usar esta fuente
para compararlas de forma directa. La uso en la sección de fundamentos del TA-IND-01 para
explicar qué es el consenso y por qué los quórums de estas bases son más pequeños que los
de un sistema BFT.

% =================== FICHA 5 ===================
\ficha{Ficha 5: Desempeño de un clúster Cassandra replicado en AWS}

\campo{Referencia (IEEE)}{\fullcite{gorbenko2021performance}}
\campo{Identificador}{DOI: \href{https://doi.org/10.32620/reks.2021.4.13}{10.32620/reks.2021.4.13}}
\campo{Tipo de fuente}{Artículo de revista de acceso abierto indexada en Scopus y DOAJ
(Radioelectronic and Computer Systems)}

\parte{Resumen}
Es un reporte experimental sobre el rendimiento de un clúster de Cassandra con tres
réplicas montado en Amazon AWS. Los autores señalan que CAP y PACELC describen la
relación entre consistencia y latencia solo de forma cualitativa, y quieren ponerle números.
Probaron cuatro formas de ubicar al cliente y a las réplicas, en una sola región o repartidas
entre regiones, usando YCSB con 100 a 1000 hilos y los niveles ONE, QUORUM y ALL.
También proponen una notación para describir esos despliegues. Con todo en una región,
la consistencia fuerte cuesta poco; cuando las réplicas están en regiones distintas, la
latencia se dispara.

\parte{Dato concreto}
En el escenario (a), con el cliente y las tres réplicas en la misma región, el throughput
de lectura bajó de 16\,830~ops/s con ONE a 13\,735~ops/s con ALL, cerca de un 18\,\% menos.
En el escenario (c), con dos réplicas en otras regiones, la lectura con ONE llegó a
15\,160~ops/s y con ALL apenas a 355~ops/s.

\campo{Ubicación}{Sección 2.2 y Tabla 2, p.~160.}

\parte{Nota personal}
Es la única de mis fuentes que estudia solo Cassandra, y además prueba QUORUM, un nivel
intermedio que las fichas 1 y 2 no incluyen. Me sirvió ver que con ONE el despliegue
multirregión rinde apenas un 13\,\% menos que el local, porque el cliente siempre manda las
peticiones al nodo más cercano. Voy a usar la Tabla 2 en la sección del TA-IND-01
dedicada a Cassandra y como referencia si repito alguna prueba con YCSB.

% =================== DECLARACIÓN DE USO DE IA ===================
\section*{Declaración de uso de inteligencia artificial}

% REVISAR: esta declaración debe describir con exactitud lo que hizo el autor.
Usé Claude (Anthropic) para escribir el código fuente LaTeX de este documento y el archivo
\texttt{referencias.bib}, y para comprobar que los metadatos de cada referencia (autores,
año, título, revista, volumen, número y páginas) coincidieran con los registros públicos de
cada artículo. Verifiqué el resultado abriendo cada DOI y comparando los datos con la
página del editor. La búsqueda y lectura de las fuentes, los resúmenes, los datos concretos
y las notas personales son de mi autoría.

\clearpage
\printbibliography[title={Referencias}]

\end{document}
